\section{Further universal four-input guarantees}\label{app:four-refinements}

This appendix records additional universal four-input divisibility guarantees
beyond the source pool used to organize the factorial-through-six proof.  Each
statement below is a separate existence theorem for a legal read-once
expression using only $+,-,\times$ and every input exactly once.

\subsection{Two auxiliary three-even blocks}\label{app:three-even-new}

The following four-input refinements use two convenient three-even strengthenings.

\begin{lemma}[Three even inputs give $20$]\label{lem:three-even20}
Any three even integers admit a read-once expression divisible by $20$.
\end{lemma}
\begin{proof}
Let the inputs be $a,b,c$.  If one is divisible by $5$, their product is
already divisible by $40$, hence by $20$.  Otherwise the three
nonzero residues modulo $5$ occupy the two signed classes
$\{\pm1\},\{\pm2\}$.  Two inputs, say $a,b$, lie in one class, so one of
$a+b,a-b$ is divisible by $5$.  It is also even, and the remaining input $c$
is even; hence $(a\pm b)c$ is divisible by $20$.
\end{proof}

\begin{lemma}[Three even inputs give $32$]\label{lem:three-even32}
Any three even integers admit a read-once expression divisible by $32$.
\end{lemma}
\begin{proof}
Let the inputs be $a,b,c$.  If $32\mid abc$, use the product.  Otherwise no
input is divisible by $8$, and two inputs cannot both be divisible by $4$.
If exactly one input, say $c$, is divisible by $4$, then $a,b\equiv2\pmod4$.
Modulo $8$ they lie in the single signed class $\{\pm2\}$, so one of
$a+b,a-b$ is divisible by $8$; multiplying by $c$ gives a multiple of $32$.

It remains that none of the three inputs is divisible by $4$.  Then all are
$2\pmod4$, and modulo $16$ they occupy only the two signed classes
$\{\pm2\}$ and $\{\pm6\}$.  Two, say $a,b$, occupy one class, so one of
$a+b,a-b$ is divisible by $16$.  Multiplication by the remaining even input
gives a multiple of $32$.
\end{proof}

\subsection{Signed-class preliminaries for four inputs}\label{app:signed-prelim}

For an odd modulus $m$, write $[r]_m$ for the signed residue class
$\{r,-r\}$ modulo $m$.  When all classes under discussion are nonzero, their
products are defined by
\[
 [r]_m[s]_m=[rs]_m.
\]
This notation is only a bookkeeping device for congruences; no division of
input values is introduced.  By \Cref{lem:signnorm}, leaves may be oriented
to convenient signed representatives and restored afterward without changing
divisibility.

\begin{lemma}[Small signed-class relation]\label{lem:small-signed-group}
Let $G$ be a finite abelian group with $|G|\le6$.  Any four-element subset
contains three distinct elements $X,Y,Z$ with $XY=Z$.
\end{lemma}
\begin{proof}
Write $G$ additively.  If the identity is not in the four-element set $A$,
fix $x\in A$.  The three sums $x+y$ with
$y\in A\setminus\{x\}$ are distinct.  If none equals a third, distinct member
of $A$, all three lie outside $A$, but there are at most two elements outside
$A$, a contradiction.  If the identity is in $A$, write
$A=\{0,x,y,z\}$.  The three sums $x+y,x+z,y+z$ are distinct, and the same
argument applies.
\end{proof}

\subsection{\texorpdfstring{The universal moduli $17$ and $19$}{The universal moduli 17 and 19}}\label{app:p17p19}

\begin{theorem}\label{thm:p17p19}
The universal guarantees $P(4;17)$ and $P(4;19)$ hold.
\end{theorem}
\begin{proof}
Fix $p\in\{17,19\}$.  If one input is divisible by $p$, multiply it by any
legal expression on the other three leaves.  If two inputs occupy the same
nonzero signed class modulo $p$, their suitable sum or difference is a
$p$-multiple and may be multiplied by the remaining two leaves.  We may
therefore assume that the four signed classes are nonzero and distinct.

First suppose none is the identity class $[1]_p$.  Consider the four
singleton classes together with their six pair-products.  If no pair-product
equals a third singleton class and no two disjoint pair-products are equal,
then these ten classes are all distinct: equality of pair-products sharing a
factor would force two singleton classes to be equal, while equality of a
pair-product with one of its own factors would force the other factor to be
$[1]_p$.  But the numbers of nonzero signed classes are $8$ for $p=17$ and
$9$ for $p=19$.  Thus a useful relation is forced.  A relation
$[a]_p[b]_p=[c]_p$ gives $(ab\pm c)d$, and equality of disjoint pair-products
gives $ab\pm cd$.

It remains that one class is $[1]_p$.  The cases not already having one of
the preceding multiplicative relations are exactly the rows in the following
table.  These rows do not enumerate integer input quadruples.  They enumerate
the residual signed residue-class configurations modulo $p$ after the preceding
general congruence reductions.  Thus, for example, the row
$(1,2,3,4)$ for $p=17$ means that, after sign normalization and relabelling,
\[
 (a,b,c,d)\equiv(1,2,3,4)\pmod {17};
\]
it does not mean that the original integers themselves are $1,2,3,4$.  The
listed construction then gives
\[
 (a+b-c)d\equiv(1+2-3)4\equiv0\pmod {17}.
\]
Every integer quadruple is covered either by one of the preceding general
relations or, after sign normalization and relabelling, by one of the rows
below.  Each row therefore represents an entire signed congruence class of
integer inputs, not a numerical test case.

\small
\begin{center}
\begin{tabular}{@{}ccl@{}}
\toprule
$p$ & representatives $(a,b,c,d)$ & expression \\
\midrule
17&(1,2,3,4)&$(a+b-c)d$\\
17&(1,2,4,5)&$(a+b)c+d$\\
17&(1,2,4,6)&$ab+c-d$\\
17&(1,2,4,7)&$a+b+c-d$\\
17&(1,2,6,7)&$a-b-c+d$\\
17&(1,3,4,8)&$a+b+c-d$\\
17&(1,3,5,8)&$a+b+c+d$\\
17&(1,4,5,8)&$(a+b-c)d$\\
17&(1,4,6,8)&$a-b-c-d$\\
17&(1,4,7,8)&$c+(a-b)d$\\
\midrule
19&(1,2,3,4)&$(a+b-c)d$\\
19&(1,2,3,5)&$(a+b-c)d$\\
19&(1,2,4,6)&$ab+c-d$\\
19&(1,2,4,7)&$a+b+c-d$\\
19&(1,2,5,6)&$a-b-c+d$\\
19&(1,2,6,8)&$ab+c-d$\\
19&(1,3,4,8)&$a+b+c-d$\\
19&(1,3,4,9)&$(a+b-c)d$\\
19&(1,3,5,9)&$a+b+c-d$\\
19&(1,3,7,9)&$(a+b)c-d$\\
19&(1,4,6,9)&$a-b-c+d$\\
19&(1,5,6,7)&$a+b+c+d$\\
19&(1,5,6,9)&$(a+b-c)d$\\
19&(1,5,8,9)&$(a+b)c+d$\\
\bottomrule
\end{tabular}
\end{center}
\normalsize
The list is exhaustive because the identity-containing four-subsets of the
signed-class multiplication table were reduced by the same two
multiplicative relations used above.  Direct substitution verifies every
row, and sign transport returns a legal expression on the original inputs.
\end{proof}

\subsection{\texorpdfstring{The universal modulus $21$}{The universal modulus 21}}\label{app:p21}

\begin{theorem}\label{thm:p21}
The universal guarantee $P(4;21)$ holds.
\end{theorem}
\begin{proof}
If one input is divisible by $7$, retain its factor $7$ and combine the
remaining three leaves into a $3$-multiple by first using
\Cref{lem:pair3} on two of them and multiplying by the third.  Otherwise the
four inputs occupy the three nonzero signed classes modulo $7$, so a
collision pair has a suitable sum or difference divisible by $7$.  The
complementary pair gives a $3$-multiple by \Cref{lem:pair3}.  The two blocks
are disjoint, hence their product is divisible by $21$.
\end{proof}

\subsection{\texorpdfstring{The universal modulus $22$}{The universal modulus 22}}\label{app:p22}

\begin{theorem}\label{thm:p22}
The universal guarantee $P(4;22)$ holds.
\end{theorem}
\begin{proof}
If one input is divisible by $11$, multiply it by an even legal expression on
the other three leaves: use their product if one of those leaves is even, and
use $(b-c)d$ if all three are odd.  If instead two nonzero inputs collide in
a signed class modulo $11$, their suitable sum or difference is an
$11$-block; the complementary two leaves can always be combined into an even
expression, using a sum or difference for equal parity and the product for
opposite parity.  Thus both situations give a $22$-multiple.

Assume henceforth that the four signed classes modulo $11$ are nonzero and
distinct.  The five nonzero signed classes form an abelian multiplication
table of order five, so \Cref{lem:small-signed-group} supplies a relation
$[a][b]=[c]$.  If the remaining leaf $d$ is even, or if $ab$ and $c$ have the
same parity, the appropriate $(ab\pm c)d$ is already even and divisible by
$11$.  A direct parity check of the possible relations shows that only the
pattern $3\OO1\EE$ can remain.

Let $u,x,y$ be the three odd leaves and $e$ the even leaf.  If
$[u][x]=[y]$ for some relabelling, then $e(ux\pm y)$ gives $22$.  Otherwise
the relation-free triples of distinct signed representatives modulo $11$ are
\[
 \{1,2,3\},\quad\{1,2,4\},\quad\{1,3,5\},\quad\{1,4,5\}.
\]
Over all relabellings of the three odd leaves, the even-valued forms
$ux\pm y$ and $u(x\pm y)$ attain every possible distinct external signed
class except one case:
\[
 \{[u],[x],[y]\}=\{[1],[2],[4]\},\qquad [e]=[3].
\]
For this last case orient the representatives as $u=1,x=2,y=4,e=3$.  Then
\[
 u-x+ye=1-2+4\cdot3=11,
\]
and the expression is even.  Thus every branch gives a multiple of $22$.
\end{proof}

\subsection{\texorpdfstring{The universal modulus $23$}{The universal modulus 23}}\label{app:p23}

\begin{theorem}\label{thm:p23}
The universal guarantee $P(4;23)$ holds.
\end{theorem}
\begin{proof}
A leaf divisible by $23$, or a collision of two nonzero signed classes
modulo $23$, immediately gives a $23$-multiple using all four leaves.  We may
therefore assume four distinct nonzero signed classes.

Suppose first that the identity class $[1]$ is absent.  Assume, toward a
contradiction, that there is no relation $[a][b]=[c]$, no equality of
disjoint pair-products, and no relation in which the product of three
classes equals the fourth.  The four singleton classes and their six
pair-products are then ten distinct classes.  Indeed, two pair-products that
share a factor cannot be equal, a pair-product cannot equal a third singleton
by assumption, and it cannot equal one of its own factors because $[1]$ is
absent.  There are eleven nonzero signed classes modulo $23$, so at most one
class is not among these ten.

Now consider the four triple-products.  They are distinct, so at least three
of them lie among the ten classes just described.  A triple-product cannot
equal a pair-product without either cancelling a common factor to obtain a
forbidden pair-to-singleton relation or forcing an identity singleton.  If a
triple-product equals a singleton different from the omitted member, then
cancelling the shared singleton shows that the other two classes form an
identity-product pair.  One such pair accounts for at most two of the four
triple-products.  Hence at least two identity-product pairs are forced.  They
must be disjoint, since the four classes are distinct and nonidentity; their
two pair-products are both $[1]$, contradicting the assumed absence of an
equality of disjoint pair-products.  Thus one of the useful relations exists,
and it gives a legal expression of the form $(ab\pm c)d$, $ab\pm cd$, or
$abc\pm d$.

It remains that one class is $[1]$.  Orient and order the signed
representatives as
\[
 a=1,
 \qquad 2\le x<y<z\le11.
\]
For two representatives put
\[
 C(r,s)=\{[rs],[r+s],[r-s]\}_{23}.
\]
If, for one of the three pairings,
\[
 C(1,x)\cap C(y,z),\quad C(1,y)\cap C(x,z),\quad
 C(1,z)\cap C(x,y)
\]
is nonempty, two legal two-leaf expressions occupy one signed class and their
suitable sum or difference gives a four-leaf multiple of $23$.  After also removing the pair-to-singleton and disjoint pair-product
relations used above, the complete residual list is
\small
\begin{center}
\begin{tabular}{@{}cccl@{}}
\toprule
$x$ & $y$ & $z$ & legal $23$-multiple \\
\midrule
2&3&9&$(a+x-y)z$\\
3&5&11&$x(a-y)-z$\\
4&5&11&$(a+x-y)z$\\
4&7&8&$x(a+y-z)$\\
4&7&9&$(a+y)x-z$\\
6&7&8&$(a+x-y)z$\\
6&8&11&$(a+x)z-y$\\
8&9&10&$(a-z+y)x$\\
9&10&11&$(a-z+y)x$\\
\bottomrule
\end{tabular}
\end{center}
\normalsize
For $x=5,7,10,11$ no residual pair $(y,z)$ remains.  The displayed table is
obtained directly from the eleven signed classes and the three displayed
$C$-intersections, so it exhausts the identity branch.  Each last-column
expression is visibly $0$ modulo $23$.  This completes the proof.
\end{proof}

\subsection{\texorpdfstring{The universal modulus $25$}{The universal modulus 25}}\label{app:p25}

\begin{theorem}\label{thm:p25}
The universal guarantee $P(4;25)$ holds.
\end{theorem}
\begin{proof}
If one input is divisible by $5$, retain that factor and use a three-leaf
$5$-multiple on the other inputs, available because $P(3;10)$ is already
proved.  Thus suppose all four leaves are units modulo $5$.  There are two
nonzero signed classes, $\{\pm1\}$ and $\{\pm2\}$.

If their multiplicities are $4+0$ or $2+2$, partition the leaves into two
same-class pairs.  Each pair has a suitable sum or difference divisible by
$5$, and the product of the two disjoint blocks is divisible by $25$.

It remains to handle a $3+1$ distribution.  Both orientations must be
considered.

\emph{Three $\{\pm1\}$ leaves and one $\{\pm2\}$ leaf.}
Orient
\[
 x=1+5\alpha,\quad y=1+5\beta,\quad z=1+5\gamma,
 \quad w=2+5\delta \pmod{25},
\]
with $\alpha,\beta,\gamma,\delta$ residues modulo $5$.  If two of
$\alpha,\beta,\gamma$ coincide, the corresponding difference is already a
$25$-block.  Otherwise put
\[
 A=\{\alpha,\beta,\gamma\},\qquad
 S=\alpha+\beta+\gamma,\qquad r=\delta-S.
\]
The following legal expressions are divisible by $25$ under the indicated
conditions:
\[
\begin{array}{c|c}
\text{condition on }r & \text{expression} \\
\hline
r=0 & xy+z-w\\
r=\alpha\ \text{(cyclically)} & x(y+z)-w\\
r=2\alpha\ \text{(cyclically)} & y+z-xw\\
r=-2\beta\ \text{(cyclically)} & (w-x)y-z
\end{array}
\]
Indeed, expansion modulo $25$ reduces these four rows respectively to
multiples of $5$ with coefficients $-r$, $\alpha-r$,
$-3\alpha-r$, and $r+2\beta$ modulo $5$.

It remains only to see that these conditions cover every $r$.  If
$0\notin A$, the three elements of $A$ are three of the four nonzero residues.
Let $t$ be the missing nonzero residue.  Then $3t$ is nonzero and different
from $t$, so $3t\in A$; doubling gives $2(3t)=t$ modulo $5$, hence
$t\in2A$.  If $0\in A$, write the two nonzero members as $u$ and $v$.  The
third nonzero direction $v$ is one of $-u,2u,-2u$.  If $v=-u$, then
$A\cup2A$ already contains $u,-u,2u,-2u$; in the other two cases the missing
nonzero direction is supplied by $2A$ or $-2A$.  Hence
\[
 \{0\}\cup A\cup2A\cup(-2A)
\]
is the full residue system modulo $5$.

\emph{Three $\{\pm2\}$ leaves and one $\{\pm1\}$ leaf.}
Orient
\[
 x=2+5\alpha,\quad y=2+5\beta,\quad z=2+5\gamma,
 \quad w=1+5\delta \pmod{25}.
\]
Again a repeated lift among $\alpha,\beta,\gamma$ gives a $25$-difference.
Assume they are distinct and use the same $S$ and $r=\delta-S$.  The family
\[
 x(y-w)-z
\]
has condition $r=2\alpha+\gamma=S+(\alpha-\beta)$ after relabelling, so as
the ordered pair $\alpha,\beta$ varies it covers every residue except $S$.
The family
\[
 z(wx-y)
\]
has condition $r=-S+2(\alpha-\beta)$ and therefore covers every residue
except $-S$.  The ordered differences of any three-element subset modulo $5$
are all four nonzero residues.  If $S\ne0$, the two families therefore cover
all $r$.

If $S=0$, both families miss only $r=0$.  The three lifts are then either
$\{0,1,-1\}$ or $\{0,2,-2\}$.  In the first case orient $\gamma=1$ and use
$z(w+x+y)$; in the second orient $\gamma=-2$ and use $z(w+xy)$.  The inner
parenthesis is a multiple of $25$ in each case.  Thus every $3+1$ branch is
covered, proving $P(4;25)$.
\end{proof}

\subsection{\texorpdfstring{The universal modulus $26$}{The universal modulus 26}}\label{app:p26}

\begin{theorem}\label{thm:p26}
The universal guarantee $P(4;26)$ holds.
\end{theorem}
\begin{proof}
If one input is divisible by $13$, multiply it by an even legal expression on
the other three leaves.  If two nonzero inputs collide in one signed class
modulo $13$, their suitable sum or difference is a $13$-block and the
complementary pair can be combined into an even expression.  Hence assume
four distinct nonzero signed classes.  By \Cref{lem:small-signed-group}, some
three distinct classes satisfy $[a][b]=[c]$.  The construction $(ab\pm c)d$
is a $26$-multiple whenever the bracket is even or $d$ is even.  A complete
parity check leaves only $2\OO2\EE$ and $3\OO1\EE$ residuals.

For $2\OO2\EE$, orient the odd leaves as $p,q$ and the even leaves as $x,y$.
After all parity-compatible pair-to-third relations have been used, exactly
six signed-class configurations remain.  The last column gives an even
expression that is $0$ modulo $13$.
\small
\begin{center}
\begin{tabular}{@{}ccc@{}}
\toprule
ordered odd representatives $(p,q)$ & ordered even representatives $(x,y)$ & expression \\
\midrule
$(1,2)$&$(3,4)$&$(p+q-x)y$\\
$(1,5)$&$(2,4)$&$x+(p+q)y$\\
$(1,5)$&$(2,6)$&$x(p+q-y)$\\
$(1,5)$&$(3,4)$&$x+(p-q)y$\\
$(1,5)$&$(3,6)$&$x(p+q-y)$\\
$(1,6)$&$(3,4)$&$p+q-x-y$\\
\bottomrule
\end{tabular}
\end{center}
\normalsize

For $3\OO1\EE$, let $u,x,y$ be odd and $e$ even.  Besides the direct
pair-to-third relations, use the even three-odd forms $xy\pm u$ and
$x(y\pm u)$.  If the signed class of one of these forms equals $[e]$, the
suitable sum or difference with $e$ is both even and divisible by $13$.
The complete signed-class reduction leaves only
\[
\begin{array}{c|c|c}
\text{odd classes}&[e]&\text{expression}\\ \hline
\{1,2,4\}&5&u-(x-e)y\\
\{1,3,5\}&4&u-x(y+e)\\
\{3,4,6\}&2&u(x-y+e)
\end{array}
\]
with representatives in the displayed order.  Their values are respectively
$13,-26,0$ modulo $13$, and all three expressions are even.  Thus every
residual case gives a $26$-multiple.
\end{proof}

\subsection{\texorpdfstring{The universal modulus $27$}{The universal modulus 27}}\label{app:p27}

\begin{theorem}\label{thm:p27}
The universal guarantee $P(4;27)$ holds.
\end{theorem}
\begin{proof}
If one input is divisible by $3$, retain its factor $3$ and use the
three-input $9$ theorem \Cref{app:three9} on the remaining leaves.  Otherwise
the four inputs are units modulo $9$, whose nonzero signed classes are
$\{\pm1\},\{\pm2\},\{\pm4\}$.  Two inputs collide in one class, producing a
$9$-block; the complementary pair gives a $3$-multiple by
\Cref{lem:pair3}.  The disjoint product is divisible by $27$.
\end{proof}

\subsection{\texorpdfstring{The universal modulus $28$}{The universal modulus 28}}\label{app:p28}

\begin{theorem}\label{thm:p28}
The universal guarantee $P(4;28)$ holds.
\end{theorem}
\begin{proof}
If an input is divisible by $7$, multiply it by an $8$-multiple from the
remaining three inputs, available by \Cref{app:three8}; the result is in fact
a multiple of $56$.  Suppose no input is divisible by $7$.  Then four inputs
occupy the three nonzero signed classes modulo $7$, so a collision exists.

If a collision pair has the same parity, its suitable sum or difference is
divisible by both $7$ and $2$, hence is a $14$-block.  The complementary pair
can be made even, giving $28$.

Assume therefore that every collision is odd-even.  For $3\OO1\EE$, the
three odd inputs give a $14$-multiple and multiplication by the even leaf
gives $28$.  For $1\OO3\EE$, if one even input is divisible by $4$, retain
that factor $4$ and use the three-input $7$ theorem on the other leaves.  If
not, all three even inputs are $2\pmod4$.  Since no even-even collision is
allowed, they occupy all three signed classes modulo $7$.  The odd input
collides with one even input, giving a $7$-block; the product of the two
remaining even inputs supplies a factor $4$.

It remains to consider $2\OO2\EE$.  If one even input is divisible by $4$,
argue as above.  Otherwise the two even inputs $x,y$ are both $2\pmod4$.
The two odd signed classes and the two even signed classes are two-element
subsets of the three nonzero signed classes modulo $7$.  If the two subsets
are equal, one cross pairing has matched classes, so the corresponding two
odd-even products occupy one signed class.  If the subsets differ, they
share one class $C$; call the other two classes $D,E$.  Direct multiplication
in the three-class table gives
\[
 C C = D E.
\]
Thus in either case one of the two cross pairings produces products in the
same signed class modulo $7$.  Their suitable sum or difference is divisible
by $7$.  Each odd-even product is $2\pmod4$, so the same sum or difference is
divisible by $4$.  Hence it is a $28$-multiple.  The all-odd and all-even
patterns cannot enter the residual branch because any collision there has
the same parity.
\end{proof}

\subsection{\texorpdfstring{The universal modulus $30$}{The universal modulus 30}}\label{app:p30}

\begin{theorem}\label{thm:p30}
The universal guarantee $P(4;30)$ holds.
\end{theorem}
\begin{proof}
If one input is divisible by $3$, retain that factor and use the three-input
$10$ theorem \Cref{app:three10}.  If one input is divisible by $5$, retain
that factor and use the three-input $12$ theorem \Cref{app:three12}; this
gives a multiple of $60$.  Hence assume all four inputs are coprime to $15$.

Modulo $5$ there are two nonzero signed classes.  If one class contains a
same-parity pair, that pair has a suitable sum or difference divisible by
both $5$ and $2$, hence a $10$-block; the complementary pair gives a
$3$-multiple.  The only residual parity pattern is therefore $2\OO2\EE$,
with each modulo-$5$ signed class containing one odd and one even input.
Orient the classes so that $a,x\equiv1\pmod5$ and
$b,y\equiv2\pmod5$, where $a,b$ are odd and $x,y$ are even.  Consider
\[
 E_1=ax+sby,
 \qquad
 E_2=ay+sbx,
 \qquad s\in\{\pm1\}.
\]
Modulo $5$, $E_1$ requires $s=+1$ while $E_2$ requires $s=-1$.
Modulo $3$, orient the four unit residues as
$\alpha,\beta,\xi,\eta\in\{\pm1\}$.  The two pairs of product terms have the
same product
\[
 (\alpha\xi)(\beta\eta)=(\alpha\eta)(\beta\xi),
\]
so the sign required to cancel the two terms modulo $3$ is the same for
$E_1$ and $E_2$.  If that common sign is $+1$, use $E_1$; if it is $-1$, use
$E_2$.  The selected expression is divisible by both $3$ and $5$, and it is
even because each product term is odd-even.  Hence it is divisible by $30$.
\end{proof}

\subsection{A complementary binary lift}\label{app:binary-lift}

\begin{lemma}[Binary complementary lift]\label{lem:binary-lift}
Suppose two legal read-once expressions satisfy
\[
 E_0\equiv M L\pmod{2M},\qquad
 E_1\equiv M(L+1)\pmod{2M}
\]
for an integer $L$.  Then one of $E_0,E_1$ is divisible by $2M$.
\end{lemma}
\begin{proof}
The consecutive integers $L$ and $L+1$ have opposite parity.
\end{proof}

\subsection{\texorpdfstring{The universal modulus $32$}{The universal modulus 32}}\label{app:p32}

\begin{theorem}\label{thm:p32}
The universal guarantee $P(4;32)$ holds.
\end{theorem}
\begin{proof}
The parity patterns other than $3\OO1\EE$ follow from the smaller
blocks already proved: four odd inputs split into an $8$-pair and a $4$-pair; $2\OO2\EE$
splits into a $4$-odd pair and an $8$-even pair; and in the $1\OO3\EE$ and
$4\EE$ patterns, three even leaves give $32$ by
\Cref{lem:three-even32}, after which the remaining leaf is multiplied in.

Let $a,b,c$ be odd and $e$ even.  If $4\mid e$, the three odd leaves give an
$8$-multiple and multiplication by $e$ gives $32$.  Hence assume
$e\equiv2\pmod4$.  If two odd leaves occupy one signed class modulo $16$,
their suitable sum or difference is a $16$-block, and multiplication by the
other odd leaf and $e$ gives $32$.  If three distinct odd classes satisfy
$[a][b]=[c]$ modulo $16$, use $e(ab\pm c)$.

The odd signed classes modulo $16$ are represented by $1,3,5,7$.  Direct
multiplication shows that the only relation-free triples are
\[
 R_A=\{1,3,7\},\qquad R_B=\{1,5,7\}.
\]
The signed class of $e$ modulo $16$ is represented by $2$ or $6$.  Orient the
leaves accordingly and write, modulo $32$,
\[
 a=r_a+16A,\quad b=r_b+16B,\quad c=r_c+16C,
 \quad e=s+16D,
\]
where $A,B,C,D\in\{0,1\}$ and $s\in\{2,6\}$.  Direct expansion gives the
following complementary pairs; every congruence is modulo $32$.
\small
\begin{center}
\begin{tabular}{@{}ccll@{}}
\toprule
odd classes & $s$ & $E_0$ & $E_1$ \\
\midrule
$(1,3,7)$&2&$c(a-b+e)\equiv16(A+B+D)$&$ce+b-a\equiv16(1+A+B+D)$\\
$(1,3,7)$&6&$c(e+b)+a\equiv16(A+B+C+D)$&$a(e+c)+b\equiv16(1+A+B+C+D)$\\
$(1,5,7)$&2&$ab+e-c\equiv16(A+B+C+D)$&$a-c(b+e)\equiv16(1+A+B+C+D)$\\
$(1,5,7)$&6&$c(a+b-e)\equiv16(A+B+D)$&$a+b+ce\equiv16(1+A+B+D)$\\
\bottomrule
\end{tabular}
\end{center}
\normalsize
Each row has the form of \Cref{lem:binary-lift} with $M=16$, so one displayed
expression is divisible by $32$.  This covers the final parity pattern.
\end{proof}

\subsection{\texorpdfstring{The universal modulus $33$}{The universal modulus 33}}\label{app:p33}

\begin{theorem}\label{thm:p33}
The universal guarantee $P(4;33)$ holds.
\end{theorem}
\begin{proof}
If one input is divisible by $11$, multiply it by a three-leaf $3$-multiple.
If two nonzero inputs collide in one signed class modulo $11$, their suitable
sum or difference is an $11$-block and the complementary pair gives a
$3$-multiple.  Thus assume four distinct nonzero signed classes modulo $11$.
Orient them to one of the five representative sets
\[
 1234,\quad1235,\quad1245,\quad1345,\quad2345.
\]
For later use, the following expressions are all $0$ modulo $11$ in their
respective rows.
\small
\begin{center}
\begin{tabular}{@{}cl@{}}
\toprule
classes $(a,b,c,d)$ & standard $11$-relations \\
\midrule
$1234$&$a(bd+c),\ b(cd-a),\ ac+bd$\\
$1235$&$a(bc+d),\ c(bd+a),\ ad+bc$\\
$1245$&$c(bd+a),\ a(cd+b),\ ab+cd$\\
$1345$&$d(bc-a),\ a(bd-c),\ ac-bd$\\
$2345$&$c(ab+d),\ d(ac+b),\ a(bd-c),\ b(cd+a),\ ad+bc$\\
\bottomrule
\end{tabular}
\end{center}
\normalsize

We now impose divisibility by $3$ on the same expression.  If at least two
leaves are $0$ modulo $3$, direct inspection of the displayed standard
relations gives one in which every surviving term is also $0$ modulo $3$.
Suppose exactly one leaf is $0$ modulo $3$.  The following table gives an
$11$-multiple on the other three leaves whenever one exists; multiplying that
three-leaf expression by the $3$-divisible leaf immediately gives $33$.
A dash marks the only four hard positions.
\scriptsize
\begin{center}
\begin{tabular}{@{}c|cccc@{}}
\toprule
classes & $a\equiv0$ & $b\equiv0$ & $c\equiv0$ & $d\equiv0$ \\
\midrule
$1234$&$bd+c$&$a-d+c$&--&$a-c+b$\\
$1235$&$b-d+c$&--&$a+bd$&$a-c+b$\\
$1245$&$b+d+c$&$a-d+c$&$a+bd$&--\\
$1345$&$bd-c$&$a-d+c$&--&$a-c+b$\\
$2345$&$bd-c$&$a+d+c$&$a-d+b$&$ac+b$\\
\bottomrule
\end{tabular}
\end{center}
\normalsize
In a hard position the remaining three residues are nonzero modulo $3$; the
two subcases in the next table cover them because two nonzero residues modulo
$3$ are either equal or opposite.
\small
\begin{center}
\begin{tabular}{@{}ccc@{}}
\toprule
classes & $3$-divisible leaf & repair \\
\midrule
$1234$&$c$&if $a\equiv b$, use $a+cd-b$; otherwise use $(a+b-c)d$\\
$1235$&$b$&if $a\equiv c$, use $(a+b-c)d$; otherwise use $(a+c)d+b$\\
$1245$&$d$&if $a\equiv c$, use $(a-c)b-d$; otherwise use $(a+c-d)b$\\
$1345$&$c$&if $a\equiv b$, use $cd-a+b$; otherwise use $(a+b-c)d$\\
\bottomrule
\end{tabular}
\end{center}
\normalsize
Every repair is $0$ modulo $11$ and, by the stated modulo-$3$ condition, also
$0$ modulo $3$.

Finally assume all four leaves are units modulo $3$.  After changing the
overall sign, normalize the modulo-$3$ residue of $a$ to $1$.  There are
exactly eight modulo-$3$ residue-sign triples for $(b,c,d)$.  Applying the
standard relations to these finite congruence possibilities leaves only the
following residual patterns; the last column gives a repair
that is simultaneously $0$ modulo $11$ and modulo $3$.
\small
\begin{center}
\begin{tabular}{@{}ccl@{}}
\toprule
classes & $(a,b,c,d)\pmod3$ & repair \\
\midrule
$1234$&$(1,-1,-1,1)$&$a+bc+d$\\
$1234$&$(1,-1,1,-1)$&$a+d-b-c$\\
$1235$&$(1,-1,1,-1)$&$a+b+c+d$\\
$1235$&$(1,1,1,1)$&$(a+c)d+b$\\
$1245$&$(1,-1,1,-1)$&$a+d-b-c$\\
$1245$&$(1,1,1,1)$&$a+d-b-c$\\
$1345$&$(1,1,-1,1)$&$a-b-c-d$\\
$1345$&$(1,-1,1,1)$&$a-b-c-d$\\
\bottomrule
\end{tabular}
\end{center}
\normalsize
The class set $2345$ has no unit residual.  The zero, one-zero, and all-unit
branches exhaust the modulo-$3$ possibilities, so $P(4;33)$ follows.
\end{proof}

\subsection{\texorpdfstring{The universal modulus $35$}{The universal modulus 35}}\label{app:p35}

\begin{theorem}\label{thm:p35}
The universal guarantee $P(4;35)$ holds.
\end{theorem}
\begin{proof}
If one input is divisible by $5$, retain that factor and apply the
three-input $7$ theorem \Cref{app:three7} to the remaining leaves.  If one
input is divisible by $7$, retain that factor and use a three-leaf
$5$-multiple on the remaining leaves.  Hence assume all four inputs are units
modulo $35$.

If two inputs occupy the same signed class modulo $35$, their suitable sum or
difference is a $35$-block.  If the four leaves can be split into two pairs
so that one pair collides modulo $5$ and the other pair collides modulo $7$
(with the roles of $5$ and $7$ interchangeable), the product of the two
suitable pair expressions is divisible by $35$.  We may also dispose of a
relation $[a][b]=[c]$ by $(ab\pm c)d$, and of an equality of disjoint
pair-products by $ab\pm cd$.

The unit signed classes modulo $35$ have representatives
\[
 1,2,3,4,6,8,9,11,12,13,16,17.
\]
After the preceding reductions, the complete residual list consists of the
following $28$ four-class sets.  As above, this is a finite classification of
signed residue classes modulo $35$, not an enumeration of integer inputs.
Each representative tuple denotes all integer quadruples that normalize to
that signed congruence configuration.  The expression in the right column is
$0$ modulo $35$; sign transport restores the original leaves.
\scriptsize
\begin{longtable}{@{}ccl@{}}
\toprule
& representatives $(a,b,c,d)$ & legal $35$-multiple \\
\midrule
\endfirsthead
\toprule
& representatives $(a,b,c,d)$ & legal $35$-multiple \\
\midrule
\endhead
1&(1,2,3,4)&$(a+b-c)d$\\
2&(1,2,3,9)&$(a+b-c)d$\\
3&(1,2,4,6)&$ab+c-d$\\
4&(1,2,4,11)&$b-(a-c)d$\\
5&(1,2,4,13)&$a+b(cd)$\\
6&(1,2,6,11)&$a-bc+d$\\
7&(1,2,8,11)&$a-b(cd)$\\
8&(1,3,6,9)&$ab+c-d$\\
9&(1,3,6,16)&$a+bc+d$\\
10&(1,3,8,16)&$a+b(cd)$\\
11&(1,3,9,13)&$a-b(cd)$\\
12&(1,3,9,16)&$a+b-cd$\\
13&(1,4,6,12)&$a-bc-d$\\
14&(1,4,8,12)&$a+b(cd)$\\
15&(1,4,11,12)&$b(a+c-d)$\\
16&(1,4,12,17)&$b(a+c)-d$\\
17&(1,6,9,17)&$a-bc-d$\\
18&(1,6,11,12)&$b(a+c-d)$\\
19&(1,6,16,17)&$b(a+c-d)$\\
20&(1,8,9,17)&$a+b(cd)$\\
21&(1,9,12,17)&$a+c-bd$\\
22&(1,9,16,17)&$a-b-cd$\\
23&(1,11,12,13)&$a-b(cd)$\\
24&(1,13,16,17)&$a-b(cd)$\\
25&(4,8,11,16)&$b-c(a-d)$\\
26&(4,11,13,16)&$b-a(c+d)$\\
27&(8,9,11,16)&$a(b-c)+d$\\
28&(9,11,13,16)&$ab-c-d$\\
\bottomrule
\end{longtable}
\normalsize
The table is exhaustive at the signed-congruence level: it is the list of
four-element subsets of the twelve displayed signed unit classes that remain
after deleting exactly the four preceding kinds of relations.  Thus no unit
configuration is omitted, and the theorem follows.
\end{proof}

\subsection{\texorpdfstring{The universal modulus $40$}{The universal modulus 40}}\label{app:p40}

\begin{theorem}\label{thm:p40}
The universal guarantee $P(4;40)$ holds.
\end{theorem}
\begin{proof}
Four odd inputs give $40$ by a modulo-$5$ collision: a collision pair is a
$10$-block and the remaining odd pair is a $4$-block; if an odd input is
itself divisible by $5$, retain its factor $5$ and use a three-odd
$8$-block.  The patterns $2\OO2\EE$ and $1\OO3\EE$ are the earlier
$40$ theorems \Cref{app:twoodd40,app:oneodd40}.  Four even inputs are handled
by choosing three even leaves for the $20$-block of
\Cref{lem:three-even20} and multiplying by the fourth.  Only $3\OO1\EE$
remains.

Let $a,b,c$ be odd and $e$ even.  If $4\mid e$, the three odd inputs give a
$10$-multiple and multiplication by $e$ gives $40$.  If an odd input is
divisible by $5$, retain its factor $5$ and use a three-input $8$-block on
the other leaves.  If $5\mid e$, then $e$ supplies a factor $10$ and the
three odd leaves supply a factor $4$.  Hence assume $e\equiv2\pmod4$ and no
input is divisible by $5$.

If two odd leaves occupy one signed class modulo $20$, their suitable sum or
difference is a $20$-block and multiplication by the third odd leaf and $e$
gives $40$.  If three distinct odd classes satisfy $[a][b]=[c]$, use
$e(ab\pm c)$.  The odd signed classes modulo $20$ are represented by
$1,3,7,9$, and direct multiplication leaves only the relation-free triples
\[
 R_A=\{1,3,9\},\qquad R_B=\{1,7,9\}.
\]
Orient the even leaf to $s=2$ or $6$ modulo $20$ and write, modulo $40$,
\[
 a=r_a+20A,\quad b=r_b+20B,\quad c=r_c+20C,
 \quad e=s+20D,
\]
where $A,B,C,D\in\{0,1\}$.  The residual lift table is
\small
\begin{center}
\begin{tabular}{@{}ccll@{}}
\toprule
odd classes & $s$ & $E_0$ & $E_1$ \\
\midrule
$(1,3,9)$&2&$c(a-b+e)\equiv20(A+B+D)$&$ce+b-a\equiv20(1+A+B+D)$\\
$(1,3,9)$&6&$ab+e-c\equiv20(A+B+C+D)$&$bc-a-e\equiv20(1+A+B+C+D)$\\
$(1,7,9)$&2&$ab+e-c\equiv20(A+B+C+D)$&$bc-a-e\equiv20(1+A+B+C+D)$\\
$(1,7,9)$&6&$c(a-b+e)\equiv20(A+B+D)$&$ce+b-a\equiv20(1+A+B+D)$\\
\bottomrule
\end{tabular}
\end{center}
\normalsize
All congruences are modulo $40$.  Each row is a binary complementary lift
with $M=20$, so one expression is divisible by $40$.
\end{proof}

